% GENERATED FILE. Edit canonical lessons, then run npm run generate:books.
% canonical-source-hash: fnv1a64:bc272aa30f134b10
% canonical-lessons: KA-C360-beteyadu, KA-C360-aledadu, KA-C360-honcu-haku, KA-C360-nuggu, KA-C360-muri

\chapter{The Hunt in the Forest}
\label{ch:ka-the-hunt-in-the-forest}
\hlchaptermodality{360}

\begin{chapteropening}
\textbf{By the end of this chapter:} \emph{I can say to hunt, to wander, to lie in wait, to barge in and to snap.}

Complete the last lesson of chapter 360: I can say to hunt, to wander, to lie in wait, to barge in and to snap.
\end{chapteropening}

\section[\texorpdfstring{\kn{ಬೇಟೆಯಾಡು} (bēṭeyāḍu)}{ಬೇಟೆಯಾಡು (bēṭeyāḍu)}]{\kn{ಬೇಟೆಯಾಡು} (bēṭeyāḍu) — to hunt}

\label{lesson:KA-C360-beteyadu}



\begin{quote}
Before the new one: say the Kannada for to parade, then the Kannada for a flag.
\end{quote}

\subsection*{You'll want to know: \kn{ಬೇಟೆಯಾಡು}}
\textbf{\kn{ಬೇಟೆಯಾಡು}} — \emph{bēṭeyāḍu} — \textquotedblleft{}to hunt\textquotedblright{}.

\begin{quote}
\kn{ಹಳೆಯ} \kn{ಕಾಲದಲ್ಲಿ} \kn{ರಾಜರು} \kn{ಕಾಡಿನಲ್ಲಿ} \kn{ಬೇಟೆಯಾಡುತ್ತಿದ್ದರು}.
\end{quote}

\begin{quote}
\emph{haḷeya kāladalli rājaru kāḍinalli bēṭeyāḍuttiddaru.}
\end{quote}

\begin{quote}
In the old days kings used to hunt in the forest.
\end{quote}

Past (I …): \kn{ಬೇಟೆಯಾಡಿದೆ}, \emph{bēṭeyāḍide}.

\kn{ಬೇಟೆ}, a hunt, and \kn{ಆಡು}, to play: to play the hunt. A hunter is a \kn{ಬೇಟೆಗಾರ}.

\begin{cousinweb}
\textbf{\kn{ಬೇಟೆಯಾಡು}} (\emph{bēṭeyāḍu}) is native Dravidian: DEDR 5527 lists Kannada \emph{bēṭe}, \textquotedblleft{}hunting, the chase\textquotedblright{}, and \emph{bēṭey āḍu}, \textquotedblleft{}to hunt\textquotedblright{}.
\end{cousinweb}

\begin{grammarlens}[title={What you've built}]
One more word for telling what happened.
\end{grammarlens}

\subsection*{Your turn}
\begin{itemize}
\raggedright
  \item \emph{Say it:} \emph{bēṭeyāḍu}
  \item \emph{Say it:} \emph{bēṭeyāḍu}, once more
  \item \emph{Say it:} \emph{haḷeya kāladalli rājaru kāḍinalli bēṭeyāḍuttiddaru.} — In the old days kings used to hunt in the forest.
  \item \emph{Recall:} say the Kannada for to parade, then the Kannada for a flag, then say \emph{bēṭeyāḍu} again
\end{itemize}

\begin{tcolorbox}[breakable,colback=teal!4,colframe=teal!35!black,title={Before you move on}]
What does \kn{ಬೇಟೆಯಾಡು} mean? (\textquotedblleft{}To hunt\textquotedblright{}.) Say it once more.
\end{tcolorbox}
\section[\texorpdfstring{\kn{ಅಲೆದಾಡು} (aledāḍu)}{ಅಲೆದಾಡು (aledāḍu)}]{\kn{ಅಲೆದಾಡು} (aledāḍu) — to wander; to roam}

\label{lesson:KA-C360-aledadu}



\begin{quote}
Before the new one: say the Kannada for a flag, then the Kannada for to hunt.
\end{quote}

\subsection*{You'll want to know: \kn{ಅಲೆದಾಡು}}
\textbf{\kn{ಅಲೆದಾಡು}} — \emph{aledāḍu} — \textquotedblleft{}to wander; to roam\textquotedblright{}.

\begin{quote}
\kn{ಕಳೆದುಹೋದ} \kn{ಹುಡುಗ} \kn{ಇಡೀ} \kn{ರಾತ್ರಿ} \kn{ಕಾಡಿನಲ್ಲಿ} \kn{ಅಲೆದಾಡಿದನು}.
\end{quote}

\begin{quote}
\emph{kaḷeduhōda huḍuga iḍī rātri kāḍinalli aledāḍidanu.}
\end{quote}

\begin{quote}
The lost boy wandered in the forest all night.
\end{quote}

Past (I …): \kn{ಅಲೆದಾಡಿದೆ}, \emph{aledāḍide}.

To move about with no fixed path. \kn{ಸುತ್ತಾಡು} is to stroll around for pleasure.

\begin{cousinweb}
\textbf{\kn{ಅಲೆದಾಡು}} (\emph{aledāḍu}) is a native Dravidian compound: \emph{ale}, \textquotedblleft{}to move about, to wander\textquotedblright{} (DEDR 240), and \emph{āḍu}, \textquotedblleft{}to move, to act\textquotedblright{} (DEDR 347).
\end{cousinweb}

\begin{grammarlens}[title={What you've built}]
One more word for telling what happened.
\end{grammarlens}

\subsection*{Your turn}
\begin{itemize}
\raggedright
  \item \emph{Say it:} \emph{aledāḍu}
  \item \emph{Say it:} \emph{aledāḍu}, once more
  \item \emph{Say it:} \emph{kaḷeduhōda huḍuga iḍī rātri kāḍinalli aledāḍidanu.} — The lost boy wandered in the forest all night.
  \item \emph{Recall:} say the Kannada for a flag, then the Kannada for to hunt, then say \emph{aledāḍu} again
\end{itemize}

\begin{tcolorbox}[breakable,colback=teal!4,colframe=teal!35!black,title={Before you move on}]
What does \kn{ಅಲೆದಾಡು} mean? (\textquotedblleft{}To wander; to roam\textquotedblright{}.) Say it once more.
\end{tcolorbox}
\section[\texorpdfstring{\kn{ಹೊಂಚು} \kn{ಹಾಕು} (hoñcu hāku)}{ಹೊಂಚು ಹಾಕು (hoñcu hāku)}]{\kn{ಹೊಂಚು} \kn{ಹಾಕು} (hoñcu hāku) — to lie in wait}

\label{lesson:KA-C360-honcu-haku}



\begin{quote}
Before the new one: say the Kannada for to hunt, then the Kannada for to wander.
\end{quote}

\subsection*{You'll want to know: \kn{ಹೊಂಚು} \kn{ಹಾಕು}}
\textbf{\kn{ಹೊಂಚು} \kn{ಹಾಕು}} — \emph{hoñcu hāku} — \textquotedblleft{}to lie in wait\textquotedblright{}.

\begin{quote}
\kn{ಹುಲಿ} \kn{ಪೊದೆಯ} \kn{ಹಿಂದೆ} \kn{ಹೊಂಚು} \kn{ಹಾಕಿತ್ತು}.
\end{quote}

\begin{quote}
\emph{huli podeya hinde hoñcu hākittu.}
\end{quote}

\begin{quote}
The tiger was lying in wait behind the bush.
\end{quote}

Past (I …): \kn{ಹೊಂಚು} \kn{ಹಾಕಿದೆ}, \emph{hoñcu hākide}.

Said of animals and of people planning an ambush.

\begin{cousinweb}
\textbf{\kn{ಹೊಂಚು} \kn{ಹಾಕು}} (\emph{hoñcu hāku}) is native Dravidian: DEDR 4596 lists Kannada \emph{hoñcu}, \textquotedblleft{}to wait for expectantly, to lurk, to lie in wait\textquotedblright{}, here with \kn{ಹಾಕು}, \textquotedblleft{}to put\textquotedblright{}.
\end{cousinweb}

\begin{grammarlens}[title={What you've built}]
One more word for telling what happened.
\end{grammarlens}

\subsection*{Your turn}
\begin{itemize}
\raggedright
  \item \emph{Say it:} \emph{hoñcu hāku}
  \item \emph{Say it:} \emph{hoñcu hāku}, once more
  \item \emph{Say it:} \emph{huli podeya hinde hoñcu hākittu.} — The tiger was lying in wait behind the bush.
  \item \emph{Recall:} say the Kannada for to hunt, then the Kannada for to wander, then say \emph{hoñcu hāku} again
\end{itemize}

\begin{tcolorbox}[breakable,colback=teal!4,colframe=teal!35!black,title={Before you move on}]
What does \kn{ಹೊಂಚು} \kn{ಹಾಕು} mean? (\textquotedblleft{}To lie in wait\textquotedblright{}.) Say it once more.
\end{tcolorbox}
\section[\texorpdfstring{\kn{ನುಗ್ಗು} (nuggu)}{ನುಗ್ಗು (nuggu)}]{\kn{ನುಗ್ಗು} (nuggu) — to barge in; to rush in}

\label{lesson:KA-C360-nuggu}



\begin{quote}
Before the new one: say the Kannada for to wander, then the Kannada for to lie in wait.
\end{quote}

\subsection*{You'll want to know: \kn{ನುಗ್ಗು}}
\textbf{\kn{ನುಗ್ಗು}} — \emph{nuggu} — \textquotedblleft{}to barge in; to rush in\textquotedblright{}.

\begin{quote}
\kn{ಆನೆ} \kn{ರಾತ್ರಿ} \kn{ಹೊಲಕ್ಕೆ} \kn{ನುಗ್ಗಿತು}.
\end{quote}

\begin{quote}
\emph{āne rātri holakke nuggitu.}
\end{quote}

\begin{quote}
An elephant barged into the field at night.
\end{quote}

Past (I …): \kn{ನುಗ್ಗಿದೆ}, \emph{nuggide}.

Where one rushes takes -\kn{ಗೆ} or -\kn{ಒಳಗೆ}: \kn{ಮನೆಯೊಳಗೆ} \kn{ನುಗ್ಗು}.

\begin{cousinweb}
\textbf{\kn{ನುಗ್ಗು}} (\emph{nuggu}) is native Dravidian: DEDR 3714 lists Kannada \emph{nuggu} beside Tamil words meaning \textquotedblleft{}to creep through, to penetrate, to get into\textquotedblright{}.
\end{cousinweb}

\begin{grammarlens}[title={What you've built}]
One more word for telling what happened.
\end{grammarlens}

\subsection*{Your turn}
\begin{itemize}
\raggedright
  \item \emph{Say it:} \emph{nuggu}
  \item \emph{Say it:} \emph{nuggu}, once more
  \item \emph{Say it:} \emph{āne rātri holakke nuggitu.} — An elephant barged into the field at night.
  \item \emph{Recall:} say the Kannada for to wander, then the Kannada for to lie in wait, then say \emph{nuggu} again
\end{itemize}

\begin{tcolorbox}[breakable,colback=teal!4,colframe=teal!35!black,title={Before you move on}]
What does \kn{ನುಗ್ಗು} mean? (\textquotedblleft{}To barge in; to rush in\textquotedblright{}.) Say it once more.
\end{tcolorbox}
\section[\texorpdfstring{\kn{ಮುರಿ} (muri)}{ಮುರಿ (muri)}]{\kn{ಮುರಿ} (muri) — to snap; to break}

\label{lesson:KA-C360-muri}



\begin{quote}
Before the new one: say the Kannada for to lie in wait, then the Kannada for to barge in.
\end{quote}

\subsection*{You'll want to know: \kn{ಮುರಿ}}
\textbf{\kn{ಮುರಿ}} — \emph{muri} — \textquotedblleft{}to snap; to break\textquotedblright{}.

\begin{quote}
\kn{ಆನೆ} \kn{ಮರದ} \kn{ದೊಡ್ಡ} \kn{ಕೊಂಬೆಯನ್ನು} \kn{ಮುರಿಯಿತು}.
\end{quote}

\begin{quote}
\emph{āne marada doḍḍa kombeyannu muriyitu.}
\end{quote}

\begin{quote}
The elephant broke a big branch of the tree.
\end{quote}

Past (I …): \kn{ಮುರಿದೆ}, \emph{muride}.

Sticks, bones and promises break: \kn{ಮಾತು} \kn{ಮುರಿ}, to break one's word. \kn{ಮುರಿತ} is a fracture.

\begin{cousinweb}
\textbf{\kn{ಮುರಿ}} (\emph{muri}) is native Dravidian: DEDR 5008 lists Kannada \emph{muṟi}, \textquotedblleft{}to break, to break off, to break down\textquotedblright{}.
\end{cousinweb}

\begin{grammarlens}[title={What you've built}]
One more word for telling what happened.
\end{grammarlens}

\subsection*{Your turn}
\begin{itemize}
\raggedright
  \item \emph{Say it:} \emph{muri}
  \item \emph{Say it:} \emph{muri}, once more
  \item \emph{Say it:} \emph{āne marada doḍḍa kombeyannu muriyitu.} — The elephant broke a big branch of the tree.
  \item \emph{Recall:} say the Kannada for to lie in wait, then the Kannada for to barge in, then say \emph{muri} again
\end{itemize}

\begin{tcolorbox}[breakable,colback=teal!4,colframe=teal!35!black,title={Before you move on}]
What does \kn{ಮುರಿ} mean? (\textquotedblleft{}To snap; to break\textquotedblright{}.) Say it once more.
\end{tcolorbox}
